\section{Conclusions and Future Work}
Despite a narrowing performance gap between proprietary and open‐source video generation systems, closed‐source solutions still lead in overall quality. Recent open‐source models such as HunyuanVideo~\cite{kong2024hunyuanvideo} and Wan2.1~\cite{wan2025} demonstrate that open frameworks can now generate realistic, high‐fidelity videos.

\textbf{Our architectural analysis reveals that:} (i) MM‐DiT and Flux‐MM‐DiT serve as the most effective backbones for modern video diffusion (ii) Flow matching has supplanted DDIM and DDPM as the preferred diffusion training objective for realism (iii) MeanFlow generates promising results that may replace Flow-matching (iv) MLLMs outperform T5 as text encoders  (v) convolutional VAEs with discriminator loss remain superior for image and video encoding (vi) Dual VAE for image and video separately showed superior results compared to one VAE for both (vii) Dual usage of 3DRoPE and 3D MM‐RoPE as positional encodings yields better temporal coherence than traditional sinusoidal embeddings (viii) LLM‐driven prompt rewriting consistently enhances generation quality.

\textbf{The current limitations are observed as follows:} (i) Substantial memory and GPU requirements that limit model scale and clip length (ii) A shortage of open‐source long‐form video datasets suited to foundational generation tasks (iii) Existing annotated datasets lack critical metadata, such as camera shot styles, camera movements, and interpersonal relationships (v) Inconsistent temporal coherence, where frame‐to‐frame continuity breaks down in longer sequences (vi) Lack of fine‐grained control over semantics in object interactions beyond coarse prompts (vii) Difficulty in modeling multiple subjects simultaneously with reference images, resulting in identity inconsistencies and unrealistic interactions (viii) Generated videos are very short between 5-16 seconds (ix) Lack of story coherent generated videos

\textbf{Solutions for the creation of long-video generation and future work:} (i) Collect long-video open-source dataset (ii) Define and annotate a hierarchical metadata schema with four key pillars: narrative segments, cinematic shot labels, character attributes (pose \& emotion), and interaction graphs (iii) Quantization and pruning for the models to overcome resources limitations (iv) Models distillation to learn from teacher models (v) Integration of prompt enhancer (v) Dividing the prompt into story narration for the coherence of the video (vi) Using multiple adapters for personalization consistency (vii) Repeating the reference image across the spatiotemporal attention.

These insights collectively chart the progress in video generation and highlight key directions for future research aimed at bridging the remaining gaps.

\label{sec:conc}